\documentclass[11pt,a4paper]{article}

\usepackage[T1]{fontenc}
\usepackage[utf8]{inputenc}
\usepackage{amsmath,amssymb}
\usepackage{booktabs}
\usepackage{listings}
\usepackage{algorithm}
\usepackage{algpseudocode}
\usepackage{geometry}
\geometry{margin=2.5cm}
\usepackage{float}
\usepackage{xcolor}
\usepackage{tikz}
\usetikzlibrary{arrows.meta,positioning}
\usepackage{hyperref}
\hypersetup{colorlinks=true,linkcolor=black,urlcolor=blue}

% hrvatski nazivi (bez babela, jer croatian jezik nije instaliran na sustavu)
\renewcommand{\abstractname}{Sažetak}
\renewcommand{\tablename}{Tablica}
\renewcommand{\figurename}{Slika}
\renewcommand{\refname}{Literatura}
\floatname{algorithm}{Algoritam}

% oznake klasa detektora (kao u prezentaciji)
\newcommand{\Dp}{\ensuremath{\mathcal{P}}}
\newcommand{\Ds}{\ensuremath{\mathcal{S}}}
\newcommand{\Dw}{\ensuremath{\mathcal{W}}}
\newcommand{\Dv}{\ensuremath{\mathcal{V}}}
\newcommand{\eP}{\ensuremath{\Diamond\mathcal{P}}}
\newcommand{\eS}{\ensuremath{\Diamond\mathcal{S}}}
\newcommand{\eW}{\ensuremath{\Diamond\mathcal{W}}}
\newcommand{\eV}{\ensuremath{\Diamond\mathcal{V}}}
\newcommand{\Crashed}{\mathrm{Crashed}}
\newcommand{\Up}{\mathrm{Up}}

\lstset{
  basicstyle=\ttfamily\footnotesize,
  keywordstyle=\bfseries,
  commentstyle=\itshape\color{gray!70!black},
  numbers=left,numberstyle=\tiny\color{gray},
  frame=single,breaklines=true,
  showstringspaces=false,
  language=Java,
  literate={č}{{\v{c}}}1 {ć}{{\'c}}1 {ž}{{\v{z}}}1 {š}{{\v{s}}}1 {đ}{{\dj}}1
           {Č}{{\v{C}}}1 {Ć}{{\'C}}1 {Ž}{{\v{Z}}}1 {Š}{{\v{S}}}1 {Đ}{{\DJ}}1
}

\title{Otkrivači grešaka u distribuiranim sustavima}
\author{Bartol Gracin \and Vid Joha \and Matija Njegač}
\date{Distribuirani procesi}

\begin{document}
\maketitle

\section{Asinkroni sustavi i otpornost na greške}

Distribuirani sustav je otporan na greške ako nastavlja ispravno raditi i kad
neki procesi zakažu. Dvije glavne paradigme u izgradnji takvih sustava su
konsenzus i atomsko emitiranje. Konsenzus omogućuje skupu procesa da postignu
zajedničku odluku koja ovisi o njihovim početnim vrijednostima, bez obzira na
kvarove. Kod atomskog emitiranja procesi pouzdano emitiraju poruke i slažu se
oko skupa isporučenih poruka i njihova redoslijeda.

Radimo u asinkronom modelu: nema gornje granice na trajanje koraka procesa ni na
kašnjenje poruke (te su veličine konačne, ali neomeđene). To uglavnom uzrokuju
nepredvidiva opterećenja sustava. Takav model je jednostavan i prenosiv jer ne
ovisi o vremenskim pretpostavkama, ali ima važnu posljedicu: konsenzus i atomsko
emitiranje ne mogu se riješiti deterministički čak ni uz jedan jedini kvar
(rezultat Fischera, Lyncha i Patersona). Uzrok je upravo to što u asinkronom
sustavu ne možemo razlikovati proces koji se srušio od procesa koji jako sporo
odgovara.

\section{Kvarovi}

Promatramo skup od $n$ procesa $Q = \{p_1, \dots, p_n\}$ povezanih pouzdanim
kanalima. Proces može zakazati samo zaustavljanjem (\emph{crash}): kažemo da se
proces $p$ srušio u trenutku $t$ ako nakon $t$ više ne izvodi nijednu akciju.
Kvarovi su trajni; srušen proces se ne oporavlja. Proces koji se nikada ne sruši
zovemo ispravnim.

Za dani tijek izvođenja $\Omega$ označavamo s $\Crashed(t,\Omega)$ skup procesa
koji su se srušili do trenutka $t$, a $\Up(t,\Omega) = Q \setminus
\Crashed(t,\Omega)$ skup onih koji do tada rade. Pretpostavljamo da je u svakom
tijeku barem jedan proces ispravan.

\section{Otkrivači kvarova}

Otkrivač kvarova $D$ je distribuirani orakl koji daje naznake o kvarovima. Svaki
proces $p_i$ ima svoj lokalni modul $D_i$, koji prati ostale procese i održava
popis onih za koje trenutno sumnja da su se srušili. Sumnja se temelji na tome
koliko dugo $p_i$ nije primio poruku od promatranog procesa.

Cijeli otkrivač je vektor $D = \langle D_{p_1}, D_{p_2}, \dots, D_{p_n} \rangle$.
Pišemo $D_p(t)$ za skup procesa za koje modul procesa $p$ sumnja u trenutku $t$;
ako je $q \in D_p(t)$, kažemo da $p$ sumnja na $q$ u trenutku $t$.

Bitno je da otkrivači mogu pogriješiti. Ispravan proces može završiti na popisu
sumnjivih, pa kasnije biti maknut kad modul uvidi grešku. Procesi se tako mogu
dodavati i uklanjati s popisa proizvoljno mnogo puta, a u istom trenutku dva
procesa mogu imati različite popise. Da bismo razlikovali koristan otkrivač od
beskorisnog, uvodimo dva svojstva: potpunost i točnost.

\section{Potpunost}

Potpunost traži da otkrivač nakon nekog vremena posumnja na sve procese koji su
se srušili.

\begin{itemize}
  \item \textbf{Jaka potpunost:} nakon nekog trenutka svaki srušeni proces trajno
  je sumnjiv \emph{svakom} ispravnom procesu.
  \[
    \forall \Omega\ \forall p \in \Crashed(\Omega)\ \forall q \in \Up(\Omega)\
    \exists t : \forall t' \ge t,\ p \in D_q(t').
  \]
  \item \textbf{Slaba potpunost:} nakon nekog trenutka svaki srušeni proces trajno
  je sumnjiv \emph{nekom} ispravnom procesu.
  \[
    \forall \Omega\ \forall p \in \Crashed(\Omega)\ \exists q \in \Up(\Omega)\
    \exists t : \forall t' \ge t,\ p \in D_q(t').
  \]
\end{itemize}

Sama potpunost nije dovoljna: otkrivač koji stalno sumnja na sve procese
zadovoljava jaku potpunost, a beskoristan je. Zato tražimo i točnost.

\section{Točnost}

Točnost ograničava greške, to jest sumnju na ispravan proces.

\begin{itemize}
  \item \textbf{Jaka točnost:} nijedan ispravan proces nikada nije sumnjiv
  nijednom ispravnom procesu.
  \[
    \forall \Omega\ \forall t\ \forall p,q \in \Up(t) : p \notin D_q(t).
  \]
  \item \textbf{Slaba točnost:} neki ispravan proces nikada nije sumnjiv nijednom
  ispravnom procesu.
  \[
    \forall \Omega\ \exists p \in \Up(\Omega)\ \forall t\ \forall q \in \Up(t) :
    p \notin D_q(t).
  \]
\end{itemize}

\subsection*{Konačna točnost}

I slabu je točnost teško postići, jer otkrivač može nakratko posumnjati na
ispravan proces i tek kasnije ispraviti grešku. Zato dodatno oslabljujemo zahtjev
i dopuštamo rane greške, uz uvjet da one prestanu.

\begin{itemize}
  \item \textbf{Konačna jaka točnost:} postoji trenutak nakon kojeg nijedan
  ispravan proces više nije sumnjiv nijednom ispravnom procesu.
  \[
    \forall \Omega\ \exists t\ \forall t' \ge t\ \forall p,q \in \Up(t') :
    p \notin D_q(t').
  \]
  \item \textbf{Konačna slaba točnost:} postoji trenutak nakon kojeg neki ispravan
  proces više nije sumnjiv nijednom ispravnom procesu.
  \[
    \forall \Omega\ \exists t\ \forall t' \ge t\ \exists p \in \Up(\Omega)\
    \forall q \in \Up(\Omega) : p \notin D_q(t').
  \]
\end{itemize}

\section{Tipovi otkrivača}

Kombiniranjem dviju vrsta potpunosti i četiri vrste točnosti dobivamo osam klasa.

\begin{table}[H]
\centering
\begin{tabular}{lllc}
\toprule
Klasa & Potpunost & Točnost & Oznaka \\
\midrule
Savršeni              & jaka  & jaka           & \Dp \\
Konačno savršeni      & jaka  & konačna jaka   & \eP \\
Jaki                  & jaka  & slaba          & \Ds \\
Konačno jaki          & jaka  & konačna slaba  & \eS \\
Slabi                 & slaba & slaba          & \Dw \\
Konačno slabi         & slaba & konačna slaba  & \eW \\
Slaba potpunost i jaka točnost                   & slaba & jaka           & \Dv \\
Slaba potpunost i konačna jaka točnost                   & slaba & konačna jaka   & \eV \\
\bottomrule
\end{tabular}
\label{tab:klase}
\end{table}

\section{Reduciranje otkrivača}

Otkrivač $D$ je reducibilan na otkrivač $E$ ako postoji distribuirani algoritam
$T_{D \to E}$ koji izlaz od $D$ pretvara u izlaz sa svojstvima od $E$. Tada je
svaki problem rješiv pomoću $E$ rješiv i pomoću $D$, pa kažemo da je $E$ slabiji
od $D$ i pišemo $D \succeq E$. Ako vrijedi i $E \succeq D$, otkrivači su
ekvivalentni, $D \equiv E$.

Iz trivijalne redukcije (svaki proces periodički prepiše izlaz svog modula) odmah
slijedi
\[
  \Dp \succeq \Dv, \qquad \Ds \succeq \Dw, \qquad
  \eP \succeq \eV, \qquad \eS \succeq \eW.
\]
Ključna je obrnuta redukcija (Algoritam~15.1 u knjizi): svaki proces periodički
šalje svoj popis sumnjivih ostalima, čime se slaba potpunost pojačava u jaku uz
očuvanje točnosti. Odatle
\[
  \Dv \succeq \Dp, \qquad \Dw \succeq \Ds, \qquad
  \eV \succeq \eP, \qquad \eW \succeq \eS,
\]
pa vrijede ekvivalencije
\[
  \Dp \equiv \Dv, \qquad \Ds \equiv \Dw, \qquad
  \eP \equiv \eV, \qquad \eS \equiv \eW.
\]

Posljedica je važna: dovoljno je riješiti problem za četiri klase s jakom
potpunošću ($\Dp, \eP, \Ds, \eS$) i automatski je riješen za preostale četiri. K
tome je $\Dp \succeq \Ds$ i $\eP \succeq \eS$ (Slika~\ref{fig:hijerarhija}), pa se
za konsenzus sve svodi na dvije netrivijalne klase: $\Ds$ i $\eS$.

\begin{figure}[h]
\centering
\begin{tikzpicture}[>=Latex, node distance=14mm and 22mm, every node/.style={font=\small}]
  \node (P) {$\Dp \equiv \Dv$};
  \node (eP) [below left=of P] {$\eP \equiv \eV$};
  \node (S) [below right=of P] {$\Ds \equiv \Dw$};
  \node (eS) [below right=of eP] {$\eS \equiv \eW$};
  \draw[->] (P) -- (eP);
  \draw[->] (P) -- (S);
  \draw[->] (eP) -- (eS);
  \draw[->] (S) -- (eS);
\end{tikzpicture}
\caption{Hijerarhija po reducibilnosti (strelica $A \to B$ znači $A \succeq B$,
tj.\ $B$ je slabiji). Gore je najjači \Dp, dolje najslabiji \eS.}
\label{fig:hijerarhija}
\end{figure}

\section{Konsenzus pomoću otkrivača}

U konsenzusu svaki proces predlaže vrijednost, a rješenje mora zadovoljiti:
konačnost (svaki ispravan proces odluči u konačnom vremenu), usuglašenost
(nikoja dva ispravna procesa ne odluče različito) i valjanost (odlučena
vrijednost je jedna od predloženih).

Kao što smo vidjeli, u čisto asinkronom sustavu konsenzus je nemoguć već uz jedan
kvar (FLP), pa je otkrivač grešaka dodatak koji tu prepreku zaobilazi. Chandra i
Toueg pokazali su kako se konsenzus rješava za svih osam klasa, a zbog
ekvivalencija iz prethodnog poglavlja dovoljno je promotriti $\Ds$ i $\eS$. Dva
rezultata određuju granicu otpornosti:

\begin{itemize}
  \item \textbf{Uz \Ds{}} konsenzus se rješava bez ograničenja na broj kvarova
  (do $n-1$ srušenih procesa). Odgovarajući algoritam (Algoritam~15.3) smo
  implementirali; opisujemo ga u poglavlju~\ref{sec:kons-s}.
  \item \textbf{Uz \eS{}} konsenzus se rješava samo ako je većina procesa
  ispravna ($f < n/2$). Algoritam~15.4 koristi rotirajućeg koordinatora: u svakoj
  rundi jedan proces prikuplja procjene većine i predlaže novu vrijednost, a
  uvjet većine sprječava da se odluke različitih rundi raziđu.
\end{itemize}

Poanta je razlika u otpornosti: jaka (trajna) točnost klase $\Ds$ dovoljna je za
bilo koji broj kvarova, dok konačna točnost klase $\eS$ traži ispravnu većinu. K
tome je $\eS$ (odnosno $\eW$) najslabiji otkrivač koji uopće rješava konsenzus,
pa u praksi težimo implementirati barem $\eP$: jači je od $\eS$, a ostvariv je u
parcijalno sinkronom sustavu (poglavlje~\ref{sec:fd}).

\section{Implementacija konsenzusa uz jaki detektor \Ds{} (Algoritam 15.3)}
\label{sec:kons-s}

Ovaj dio projekta (razred \texttt{ConsensusProcess}) rješava konsenzus uz jaki
detektor $D \in \Ds$ i podnosi bilo koji broj kvarova (do $n-1$ srušenih
procesa). Svaki proces $p$ drži vektor $V_p$ svojih procjena predloženih
vrijednosti (na početku sve $\bot$, osim $V_p[p] = v_p$) i vektor $\Delta_p$
vrijednosti koje je naučio u zadnjoj rundi. Upit ,,$q \in D_p$'' iz pseudokoda
ostvaruje se ugrađenim detektorom (heartbeat poruke ,,Alive'' i metoda
\texttt{provjeriProces}), po istoj ideji kao Algoritam~15.7.

\begin{algorithm}[h]
\caption{Konsenzus uz jaki detektor $D \in \Ds$ (procedura \textsc{propose}$(v_p)$
na procesu $p$)}
\label{alg:S}
\begin{algorithmic}[1]
\State $V_p \gets \langle \bot, \dots, \bot \rangle$;\quad $V_p[p] \gets v_p$;
  \quad $\Delta_p \gets V_p$
\Statex \textbf{Faza 1:} \textbf{za} $r_p \gets 1$ \textbf{do} $n-1$
  \State \quad pošalji $(r_p, \Delta_p, p)$ svim procesima
  \State \quad čekaj dok $\forall q$: primljeno $(r_p, \Delta_q, q)$ \textbf{ili}
  $q \in D_p$ \Comment{upit detektoru}
  \State \quad $\Delta_p \gets \langle \bot, \dots, \bot \rangle$
  \State \quad \textbf{za} $k \gets 1$ \textbf{do} $n$
  \State \qquad \textbf{ako} $V_p[k] = \bot$ \textbf{i} $\exists\,(r_p, \Delta_q, q)$
  s $\Delta_q[k] \ne \bot$
  \State \qquad\quad $V_p[k] \gets \Delta_q[k]$;\quad $\Delta_p[k] \gets \Delta_q[k]$
\Statex \textbf{Faza 2:}
  \State \quad pošalji $V_p$ svim procesima
  \State \quad čekaj dok $\forall q$: primljeno $V_q$ \textbf{ili} $q \in D_p$
  \State \quad \textbf{za} $k \gets 1$ \textbf{do} $n$: \textbf{ako} $\exists$
  primljeni $V_q$ s $V_q[k] = \bot$ \textbf{tada} $V_p[k] \gets \bot$
\Statex \textbf{Faza 3:}
  \State \quad odluči prvi element vektora $V_p$ različit od $\bot$
\end{algorithmic}
\end{algorithm}

\subsection{Opis triju faza}

\begin{description}
  \item[Faza 1 ($n-1$ rundi).] U svakoj rundi proces emitira svoju $\Delta_p$
  (vrijednosti naučene u prethodnoj rundi) i čeka poruku te runde od svakog
  procesa koji nije osumnjičen. Ako je za neko $k$ vrijednost $V_p[k]$ još $\bot$,
  a u primljenoj poruci nije, proces je preuzima. Nakon $n-1$ rundi svaki
  ispravan proces zna sve vrijednosti koje su ,,preživjele''.
  \item[Faza 2 (usklađivanje).] Proces emitira svoj cijeli vektor $V_p$ i čeka
  vektore ostalih. Ako je bilo koji primljeni vektor imao $\bot$ na mjestu $k$,
  proces i sam postavlja $V_p[k] \gets \bot$. Time se izbacuju vrijednosti oko
  kojih se svi ne slažu; pokazuje se da barem jedno mjesto ostaje različito od
  $\bot$.
  \item[Faza 3 (odluka).] Svi ispravni procesi odluče prvi element vektora $V_p$
  različit od $\bot$; taj je element jednak kod svih.
\end{description}

\subsection{Implementacija u Javi}

Glavna petlja (u \texttt{ConsensusTester}) vodi runde Faze 1, zatim Fazu 2
(oznaka poruke ,,Provjera'') i Fazu 3 (odluka). Doslovni kod:

\begin{lstlisting}
    for(int runda=1; runda < numProc; runda++){
        System.out.println("--------------------------");
        System.out.println("Započeta "+runda+". runda");

        for (int j = 0; j < numProc; j++){
            if (j != myId){
                process.sendMsg(j,"Consensus",runda,process.getDelta());
                kvar();
            }
        }

        do{
            kvar();
            for (int j = 0; j < numProc; j++){
                if (j != myId){
                    process.sendMsg(j,"Alive",runda,process.getDelta());
                }
            }
            for (int j = 0; j < numProc; j++){
                if (j != myId){
                    process.provjeriProces(j);
                }
            }
            Thread.sleep(1000);
        }while(process.kolikoSamPorukaPrimioURundi(runda)+process.kolikoJeProcesaOsumnjiceno()<numProc-1);

        System.out.println("--------------------------");
        System.out.println("Dobio sam poruke od svih ili su mi postali sumnjivi");

        process.ispisiSvePorukeURundi(runda);

        process.obradiPoruke(runda);

        process.novaRunda();

    }
    for (int j = 0; j < numProc; j++){
            if (j != myId){
                process.sendMsg(j,"Provjera",numProc,process.getVector());
                kvar();
            }
        }
    do{
            kvar();
            for (int j = 0; j < numProc; j++){
                if (j != myId){
                    process.sendMsg(j,"Alive",numProc,process.getVector());
                }
            }
            for (int j = 0; j < numProc; j++){
                if (j != myId){
                    process.provjeriProces(j);
                }
            }
            Thread.sleep(1000);
        }while(process.kolikoSamZadnjihPorukaPrimio()+process.kolikoJeProcesaOsumnjiceno()<numProc-1);
    System.out.println("Završni vektor prije provjere je "+process.V);
    for (int j = 0; j < numProc; j++){
            process.provjeriZadnjePoruke(j);
        }
    System.out.println("Završni vektor nakon provjere je "+process.V);
    for(int i = 0; i < process.V.size() ; i++){
        if(process.V.get(i)!=null){
            System.out.println("Svi procesi se slažu oko vrijednosti "+process.V.get(i));
            break;
        }
    }
\end{lstlisting}

Spajanje vrijednosti u Fazi 1 (metoda \texttt{obradiPoruke}):

\begin{lstlisting}
    public void obradiPoruke(int runda){
        int numProc=Delta.size();

        for (int i = 0; i < numProc; i++) {
            Delta.set(i, null);
        }
        System.out.println("Poruke su");
        System.out.println(Poruke.get(runda-1));
        for (int i = 0; i < numProc; i++) {
            if(V.get(i)==null){
                for(int j = 0; j < Poruke.get(runda-1).size(); j++){
                    if(Poruke.get(runda-1).get(j)!=null){
                        System.out.println("Idem obraditi poruku "+Poruke.get(runda-1).get(j).getMessage()+" od procesa "+Poruke.get(runda-1).get(j).getSrcId());
                        String[] tokeni = Poruke.get(runda-1).get(j).getMessage().trim().split("\\s+");
                        Integer[] rezultat = new Integer[tokeni.length];
                        rezultat[i]=tokeni[i+1].equals("null") ? null : Integer.parseInt(tokeni[i+1]);
                        System.out.println("rezultat[i]: "+rezultat[i]);
                        if(rezultat[i]!=null){
                            V.set(i,rezultat[i]);
                            Delta.set(i,rezultat[i]);
                        }


                    }
                }
            }
            else if(V.get(i)!=null){
                        System.out.println("Vrijednost od V na indexu "+i+" nije null nego "+V.get(i));
                    }
        }
        System.out.println("Vektor V je ");
        System.out.println(V);
        System.out.println("Vektor Delta je ");
        System.out.println(Delta);
    }
\end{lstlisting}

Izbacivanje nepotvrđenih vrijednosti u Fazi 2 (metoda
\texttt{provjeriZadnjePoruke}):

\begin{lstlisting}
    public void provjeriZadnjePoruke(int proces){
        for(int i = 0; i<ZadnjePoruke.size() ; i++){
            if(ZadnjePoruke.get(i).getSrcId()==proces){
                System.out.println("Idem obraditi zadnju poruku "+ZadnjePoruke.get(i).getMessage()+" od procesa "+ZadnjePoruke.get(i).getSrcId());
                String[] tokeni = ZadnjePoruke.get(i).getMessage().trim().split("\\s+");
                Integer[] rezultat = new Integer[tokeni.length];
                rezultat[i]=tokeni[i+1].equals("null") ? null : Integer.parseInt(tokeni[i+1]);

                if(rezultat[i]==null){
                    V.set(i, null);
                }
            }
        }
    }
\end{lstlisting}

\subsection{Dokaz ispravnosti}

\paragraph{Otpornost.} Algoritam podnosi do $n-1$ kvarova jer nijedna faza ne
traži većinu: proces čeka poruku od svakog neosumnjičenog procesa, a osumnjičene
preskače.

\paragraph{Dogovor.} Slaba točnost detektora $\Ds$ jamči da postoji ispravan
proces $p^\ast$ kojeg nitko nikada ne osumnjiči. Zato ga svi u svakoj rundi
čekaju, pa se njegova vrijednost proširi u sve vektore i preživi Fazu 2. Svi
ispravni procesi tako u Fazi 3 vide isti prvi element različit od $\bot$ i
odluče jednako.

\paragraph{Završetak.} Jaka potpunost detektora jamči da se srušeni procesi na
kraju osumnjiče, pa čekanje u fazama nikada ne blokira zauvijek.

\section{Implementacija otkrivača \eP{} (Algoritam 15.7)}
\label{sec:fd}

Savršen otkrivač \Dp{} ne možemo implementirati u parcijalnoj sinkroniji: dok ne
znamo gornju granicu kašnjenja, svaki fiksni istek vremena može biti prekratak i
osumnjičiti spor, ali ispravan proces. Najbolje što postižemo jest da takve
greške \emph{na kraju} prestanu, a to je upravo \eP.

Parcijalna sinkronost znači da je sustav u početku asinkron, ali nakon nekog
nepoznatog trenutka postane sinkron, s omeđenim kašnjenjima. Svaki proces $p$
drži zadani istek vremena za svaki drugi proces i mjeri vrijeme lokalnim satom
koji broji korake (otkucaje). Koriste se varijable: $Output_p$ (popis sumnjivih,
na početku prazan) i $\Delta_p(q)$ (trajanje isteka vremena procesa $p$ za $q$).

\begin{algorithm}[h]
\caption{Timeout implementacija $D \in \eP$ u modelu parcijalne sinkronosti
(svaki proces $p$)}
\label{alg:eP}
\begin{algorithmic}[1]
\State $Output_p \gets \emptyset$
\ForAll{$q \in Q$}\ $\Delta_p(q) \gets$ zadani istek vremena \EndFor
\Statex \textbf{cobegin}
\Statex \textbf{Zadatak 1:} periodički
  \State \quad pošalji ,,$p$ je živ'' svima
\Statex \textbf{Zadatak 2:} periodički, za svaki $q \in Q$
  \State \quad \textbf{ako} $q \notin Output_p$ \textbf{i} $p$ nije primio
  ,,$q$ je živ'' u zadnjih $\Delta_p(q)$ otkucaja
  \State \qquad $Output_p \gets Output_p \cup \{q\}$
\Statex \textbf{Zadatak 3:} kad primiš ,,$q$ je živ''
  \State \quad \textbf{ako} $q \in Output_p$
  \State \qquad $Output_p \gets Output_p \setminus \{q\}$
  \Comment{$p$ se kaje}
  \State \qquad $\Delta_p(q) \gets \Delta_p(q) + 1$
  \Comment{i poveća istek vremena}
\Statex \textbf{coend}
\end{algorithmic}
\end{algorithm}

\subsection{Opis triju zadataka}

\begin{description}
  \item[Zadatak 1 (otkucaji srca).] Svaki proces svima periodički šalje poruku
  ,,živ sam''. Kad se proces sruši, prestane ih slati.
  \item[Zadatak 2 (istek vremena).] Ako $p$ od $q$ nije čuo ništa u zadnjih
  $\Delta_p(q)$ otkucaja, doda $q$ među sumnjive.
  \item[Zadatak 3 (kajanje).] Ako $p$ primi poruku od procesa kojeg sumnjiči, zna
  da je istek vremena bio preuranjen: makne $q$ iz sumnjivih i \emph{poveća}
  $\Delta_p(q)$, da idući put bude strpljiviji.
\end{description}

\subsection{Implementacija u Javi}

Otkrivač smo izdvojili u zaseban razred \texttt{FailureDetector} unutar okvira za
razmjenu poruka. Stanje odgovara varijablama iz pseudokoda: \texttt{Output} je
skup sumnjivih, \texttt{delta[q]} istek vremena za $q$, \texttt{lastAlive[q]}
otkucaj zadnje poruke od $q$, a \texttt{clock} lokalni sat.

\begin{lstlisting}
// Zadatak 1: posalji "ziv sam" svima
public void sendHeartbeats() {
    if (crashed) return;
    broadcastMsg("alive", (int) clock);
}

// Zadatak 2: istek vremena -> sumnja
public synchronized void checkTimeouts() {
    for (int q = 0; q < N; q++) {
        if (q == myId) continue;
        if (!Output.contains(q) && (clock - lastAlive[q]) > delta[q]) {
            Output.add(q);
        }
    }
}

// Zadatak 3: kad primim "ziv sam" od q -> kajanje + veci timeout
public synchronized void handleMsg(Msg m, int q, String tag) {
    if (tag.equals("alive")) {
        lastAlive[q] = clock;
        if (Output.contains(q)) {
            Output.remove(q);
            delta[q] = delta[q] + 1;
        }
    }
}
\end{lstlisting}

Zadaci 1 i 2 izvode se u glavnoj petlji, gdje jedan prolaz odgovara jednom
otkucaju (\texttt{clock++}). Zadatak 3 izvodi se asinkrono, u dretvama koje
slušaju poruke, pa su metode koje dijele stanje označene s \texttt{synchronized}.

\subsection{Dokaz ispravnosti}

\paragraph{Jaka potpunost.} Srušeni proces trajno prestaje slati otkucaje.
Razlika \texttt{clock - lastAlive[q]} tada raste bez granice i prije ili kasnije
premaši $\Delta_p(q)$, pa $q$ ulazi u $Output_p$. Iz $Output_p$ se izlazi samo
primitkom nove poruke od $q$, koje više nema, pa sumnja ostaje trajna.

\paragraph{Konačna jaka točnost.} Nakon svake pogrešne sumnje na ispravan proces
$q$, pri sljedećoj poruci od $q$ istek $\Delta_p(q)$ poraste za jedan. Kad sustav
postane sinkron, kašnjenja imaju konačnu gornju granicu; $\Delta_p(q)$ nakon
konačno mnogo grešaka premaši tu granicu i pogrešnih sumnji više nema.

\subsection{Testiranje}

Otkrivač smo testirali s tri procesa i scenarijima koji izdvajaju svako svojstvo
(Tablica~\ref{tab:test}).

\begin{table}[h]
\centering
\begin{tabular}{lp{9cm}}
\toprule
Scenarij & Opažanje \\
\midrule
Kvar procesa $p_2$ u otkucaju 10 &
$p_0$ i $p_1$ počnu sumnjati na $p_2$ oko otkucaja 15 i zadrže ga u $Output$ do
kraja; ispravne procese nikada ne sumnjiče. Potvrđuje jaku potpunost. \\
\addlinespace
Spor proces $p_2$ (otkucaj srca svakih 8 otkucaja) &
$p_0$ i $p_1$ ga isprva više puta pogrešno osumnjiče i pokaju se, a $\Delta$ raste
$5 \to 6 \to 7 \to 8$; nakon toga sumnje prestaju. Potvrđuje konačnu jaku
točnost. \\
\bottomrule
\end{tabular}
\caption{Rezultati testiranja otkrivača \eP.}
\label{tab:test}
\end{table}

Drugi scenarij dobro pokazuje zašto je ovo \eP, a ne \Dp. Otkrivač pogriješi dok
ne nauči pravi ritam sporog procesa, ali čim $\Delta$ dostigne stvarni razmak
otkucaja, greške staju. Savršen otkrivač tu ne bi pogriješio nijednom, ali u
parcijalnoj sinkroniji nije ostvariv.

\begin{thebibliography}{9}
\bibitem{ks}Kshemkalyani A.D., Singal M. Distributed Computing: Principles, Algorithms and Systems. Cambridge University Press, 2011.
\end{thebibliography}

\end{document}
